\documentclass[11pt]{article}
\input{../../shared/project_style.tex}
\rhead{Basics 34 Textbook Draft}

\begin{document}
\DocTitle{Basics 34: Parallel Wave Number and Field-Line Propagation}{V - Plasma waves and Alfvén physics}{Directional derivatives, toroidal mode numbers, rotational transform, rational surfaces, and sign conventions}

\begin{overviewbox}
\textbf{Textbook draft, version 1.0.} Alfvén waves propagate along magnetic field lines, so the physically relevant wave number is the component parallel to the equilibrium field. This chapter derives the toroidal expression connecting $k_\parallel$ to poloidal and toroidal mode numbers and rotational transform. This chapter is designed for a reader with no prior plasma-physics training. It distinguishes exact identities from physical approximations and connects the central equations to later simulation modules.
\end{overviewbox}

\ProjectTOC

\section{Learning objectives}
After completing this note, the reader should be able to:
\begin{itemize}[leftmargin=1.6em]
\item Define the derivative parallel to a magnetic field.
\item Relate local phase variation to parallel wave number.
\item Derive the large-aspect-ratio toroidal form of $k_\parallel$.
\item Interpret rotational transform and safety factor.
\item Locate rational or resonant magnetic surfaces.
\item Track sign conventions consistently.
\end{itemize}

\section{Prerequisites and reading strategy}
\begin{itemize}[leftmargin=1.6em]
\item Basics 2, 7, 31, and 32.
\end{itemize}
On a first reading, emphasize physical meaning and dimensional checks. On a second reading, reproduce the derivations. Attempt the computational laboratory only after the limiting cases are understood.

\section{Parallel derivative}

Let
\begin{equation}
\mathbf b=\frac{\mathbf B}{B}
\end{equation}
be the magnetic-field unit vector. The derivative along a field line is
\begin{equation}
\nabla_\parallel=\mathbf b\cdot\nabla.
\end{equation}
For a local phase $S(\mathbf x)$ with $\mathbf k=\nabla S$,
\begin{equation}
k_\parallel=\mathbf b\cdot\mathbf k=\mathbf b\cdot\nabla S.
\end{equation}
The shear-Alfvén relation is locally $\omega=\pm k_\parallel v_A$.


\section{Straight-field interpretation}

For $\mathbf B=B_0\hat{\mathbf z}$ and phase
\begin{equation}
S=k_xx+k_yy+k_zz-\omega t,
\end{equation}
the parallel wave number is simply $k_\parallel=k_z$. Components perpendicular to $\mathbf B$ shape the transverse structure but do not enter the elementary tension-wave frequency directly.


\section{Magnetic field lines on a torus}

Let $\theta$ be poloidal angle and $\phi$ toroidal angle. Rotational transform is
\begin{equation}
\iota=\frac{d\theta}{d\phi}
\end{equation}
along a magnetic field line, while tokamak safety factor is commonly $q=1/\iota$. On a rational surface $\iota=n/m$, a field line closes after $m$ toroidal or $n$ poloidal turns under the stated convention.


\section{Mode phase and toroidal $k_\parallel$}

Use the harmonic convention
\begin{equation}
\tilde f(s,\theta,\phi,t)
=\hat f(s)e^{i(m\theta-n\phi-\omega t)}.
\end{equation}
Along a field line,
\begin{equation}
\frac{d}{d\phi}(m\theta-n\phi)=m\iota-n.
\end{equation}
If toroidal arc length is approximately $R_0d\phi$,
\begin{equation}
\boxed{k_\parallel(s)\simeq\frac{m\iota(s)-n}{R_0}}.
\end{equation}
The opposite phase convention or reversed angle orientation changes the sign. Frequencies based on $k_\parallel^2$ are unchanged, but propagation direction and resonance signs are not.


\section{Rational surfaces and local continuum zeros}

At
\begin{equation}
m\iota(s_r)-n=0,
\end{equation}
the selected harmonic has $k_\parallel=0$. The elementary local Alfvén frequency therefore approaches zero:
\begin{equation}
\omega_{A,m}(s)=|k_{\parallel,m}(s)|v_A(s).
\end{equation}
In a complete toroidal system, pressure, compressibility, geometry, and coupling modify the spectrum, but rational surfaces remain central landmarks.


\section{Magnetic shear}

Magnetic shear measures radial variation of field-line pitch. One common tokamak form is
\begin{equation}
\hat s=\frac{r}{q}\frac{dq}{dr}.
\end{equation}
Because $k_\parallel$ depends on $\iota(s)$, shear changes the radial slope of continuum branches and the width of resonant layers. Low shear can produce broad near-resonant regions; strong shear changes $k_\parallel$ rapidly with radius.


\section{Three-dimensional stellarator extension}

A stellarator equilibrium has many Fourier harmonics in magnetic geometry. A perturbation component labeled $(m,n)$ can couple to $(m+\mu,n+\nu N_{\rm fp})$, where $(\mu,\nu)$ label equilibrium harmonics and $N_{\rm fp}$ is field periodicity. There is no single isolated $k_\parallel$ branch once coupling is strong; instead one solves a matrix or differential eigenproblem.


\section{Connection to particle resonance}

Wave-particle interaction depends on phase coherence along orbits. A schematic resonance is
\begin{equation}
\omega-n\omega_\phi-m\omega_\theta-\ell\omega_b-\omega_d\simeq0.
\end{equation}
Consistent mode-number and angle conventions are essential when connecting MHD eigenfunctions to orbit frequencies. A sign error can exchange co- and counter-passing resonances.


\section{Computational laboratory}
Define analytic $\iota(s)$ profiles and plot $k_{\parallel,m}(s)$ for several $(m,n)$ pairs. Locate rational surfaces, vary magnetic shear, and compare sign conventions. Trace a field line numerically and verify the phase accumulation predicted by $m\iota-n$.

\section{Summary}
\begin{itemize}[leftmargin=1.6em]
\item Parallel wave number measures phase variation along equilibrium magnetic field lines.
\item Toroidal mode labels and rotational transform determine $k_\parallel$.
\item Rational surfaces occur where a selected harmonic has zero parallel wave number.
\item Magnetic shear controls radial variation of the local parallel phase.
\item Stellarator geometry couples many harmonics, so conventions must remain consistent from fields to particle resonances.
\end{itemize}

\SymbolsAcronymsSection
\begin{symtable}
$\mathbf b$ & field unit vector & $\mathbf B/B$. \\
$k_\parallel$ & parallel wave number & $\mathbf b\cdot\nabla S$. \\
$m,n$ & mode numbers & Poloidal and toroidal indices in this note. \\
$\iota$ & rotational transform & $d\theta/d\phi$ along a field line. \\
$q$ & safety factor & Commonly $1/\iota$. \\
$\hat s$ & magnetic shear & Radial change of field-line pitch. \\
\end{symtable}

\section{Exercises}
\begin{enumerate}[leftmargin=1.8em]
\item Derive $k_\parallel=\mathbf b\cdot\nabla S$.
\item Derive the large-aspect-ratio expression from the toroidal harmonic phase.
\item Locate the rational surface for a specified $\iota(s)$, $m$, and $n$.
\item Explain why changing the Fourier convention changes propagation signs but not $\omega^2$.
\item Describe how three-dimensional equilibrium harmonics expand the coupling network.
\end{enumerate}

\section{References}
\begin{thebibliography}{99}
\bibitem{ref1} G. B. Arfken, H. J. Weber, and F. E. Harris, \emph{Mathematical Methods for Physicists}, 7th ed., Academic Press (2013).
\bibitem{ref2} G. Strang, \emph{Linear Algebra and Its Applications}, 4th ed., Brooks/Cole (2006).
\bibitem{ref3} W. A. Strauss, \emph{Partial Differential Equations: An Introduction}, 2nd ed., Wiley (2007).
\bibitem{ref4} F. F. Chen, \emph{Introduction to Plasma Physics and Controlled Fusion}, 3rd ed., Springer (2016).
\bibitem{ref5} R. J. Goldston and P. H. Rutherford, \emph{Introduction to Plasma Physics}, CRC Press (1995).
\bibitem{ref6} P. M. Bellan, \emph{Fundamentals of Plasma Physics}, Cambridge University Press (2006).
\bibitem{ref7} J. P. Freidberg, \emph{Ideal MHD}, Cambridge University Press (2014).
\bibitem{ref8} J. P. Goedbloed, R. Keppens, and S. Poedts, \emph{Advanced Magnetohydrodynamics}, Cambridge University Press (2010).
\bibitem{ref9} H. Goedbloed and S. Poedts, \emph{Principles of Magnetohydrodynamics}, Cambridge University Press (2004).
\bibitem{ref10} T. H. Stix, \emph{Waves in Plasmas}, American Institute of Physics (1992).
\bibitem{ref11} D. A. Gurnett and A. Bhattacharjee, \emph{Introduction to Plasma Physics}, 2nd ed., Cambridge University Press (2017).
\end{thebibliography}

\end{document}
